\begin{abstract}
  The NP-hard \textsc{Odd Cycle Transversal} problem asks for a minimum vertex set whose removal from an undirected input graph~$G$ breaks all odd cycles, and thereby yields a bipartite graph. The problem is well-known to be fixed-parameter tractable when parameterized by the size~$k$ of the desired solution. It also admits a randomized kernelization of polynomial size, using the celebrated matroid toolkit by Kratsch and Wahlstr\"{o}m. The kernelization guarantees a reduction in the total \emph{size} of an input graph, but does not guarantee any decrease in the size of the solution to be sought; the latter governs the size of the search space for FPT algorithms parameterized by~$k$. We investigate under which conditions an efficient algorithm can detect one or more vertices that belong to an optimal solution to \textsc{Odd Cycle Transversal}. By drawing inspiration from the popular \emph{crown reduction} rule for \textsc{Vertex Cover}, and the notion of \emph{antler decompositions} that was recently proposed for \textsc{Feedback Vertex Set}, we introduce a graph decomposition called \emph{tight odd cycle cut} that can be used to certify that a vertex set is part of an optimal odd cycle transversal. While it is NP-hard to compute such a graph decomposition, we develop parameterized algorithms to find a set of at least~$k$ vertices that belong to an optimal odd cycle transversal when the input contains a tight odd cycle cut certifying the membership of~$k$ vertices in an optimal solution. The resulting algorithm formalizes when the search space for the solution-size parameterization of \textsc{Odd Cycle Transversal} can be reduced by preprocessing. To obtain our results, we develop a graph reduction step that can be used to simplify the graph to the point that the odd cycle cut can be detected via color coding.
\end{abstract}
  